% =====================================================================
\chapter{Estructura del proyecto y control de versiones}\label{cap:proyecto}

\section{Repositorios}
\begin{itemize}
\item \textbf{Brazo-de-4-grados-de-libertad} (GitHub, \codigo{GabAndGom}): repositorio original del equipo, con el material de los semestres anteriores (2026\_I) y del actual.
\item \textbf{Brazo-4GDL-gemelo-digital} (GitHub, \codigo{GabAndGom}): derivado del anterior con todo su historial, agrega el gemelo digital. Publica la página en \url{https://brazo-4gdl-gemelo-digital.pages.dev}.
\end{itemize}
Regla del equipo: en el repositorio original no se modifica el material de la gestión I/2026; el trabajo nuevo va en \codigo{semestres/2026\_II/}.

\section{Árbol del repositorio}
\begin{lstlisting}
Brazo-4GDL-gemelo-digital/
|-- index.html, web/            pagina publicada: estilos, JS, imagenes
|   `-- js/  cinematica.js  modelo3d.js  conexion-plc.js  app.js
|-- scripts/build-catalog.js    genera el catalogo de archivos de la pagina
|-- config/                     titulos, semestres, videos (JSON)
|-- .github/workflows/pages.yml publicacion en GitHub Pages
|-- global/                     material que no depende del semestre
|   |-- componentes/            imagenes de cada componente
|   |-- datasheets/             PLC, drivers, sensores, reles
|   |-- fotos-planta/           fotos del brazo y del tablero
|   |-- manuales/               manual del fabricante, referencia ladder, MAP_SERV
|   |-- modelo-3d/              brazo_xk-js1-1.glb y su generador
|   |-- normativas/             ISA-5.1, IEC 60617
|   |-- plantillas/             plantilla LaTeX de informes
|   `-- programas-plc/          programa del fabricante, libreria MAP_SERV
`-- semestres/
    |-- 2026_I/                 gestion anterior (no se modifica)
    `-- 2026_II/
        |-- guias/              guias de laboratorio 1, 2 y 3
        |-- informes/           informe 1 (auditoria y requerimientos)
        |-- diagramas/          diagrama de bloques fenomenologico
        |-- planos/electrico/, planos/pid/   planos .dwg (historico)
        |-- base-de-datos/      mockups, modelo E-R, scripts SQL, interfaz
        |-- modelo-3d/          levantamiento, cinematica, informe y figuras
        |-- programas-plc/      programa original de secuencias (.awl, .mwp)
        |-- gemelo-digital/
        |   |-- app/            app unificada (Linux y Windows)
        |   |-- plc/            brazo_gemelo.awl y .mwp
        |   `-- manual/         manuales, arquitectura, guia y este manual
        |-- hmi/, pruebas/, evidencias/   (README con lo que va en cada una)
\end{lstlisting}

\begin{center}\small
\begin{tabularx}{\linewidth}{>{\raggedright}p{6.4cm} >{\raggedright\arraybackslash}X}
\toprule
\textbf{Carpeta pedida en la guía 3} & \textbf{Equivalente en el repositorio}\\
\midrule
\codigo{/documentacion/planos\_isa/} & \codigo{semestres/2026\_II/planos/} (eléctrico y P\&ID)\\
\codigo{/documentacion/manuales\_inversa/} & \codigo{semestres/2026\_II/informes/}, \codigo{modelo-3d/levantamiento/}, \codigo{global/datasheets/}\\
\codigo{/database/scripts/} & \codigo{semestres/2026\_II/base-de-datos/scripts/}\\
\bottomrule
\end{tabularx}
\end{center}

\section{Archivos fuera del repositorio}
\begin{center}\small
\begin{tabularx}{\linewidth}{>{\raggedright}p{5cm} >{\raggedright\arraybackslash}X}
\toprule
\textbf{Carpeta (en \codigo{ETN 1034/})} & \textbf{Contenido}\\
\midrule
\codigo{gemelo\_brazo\_app/} y \codigo{.zip} & app exportada lista para usar, con los datos propios de la PC: usuarios, certificado, calibración, registro\\
\codigo{respaldo\_vm/} & copias de los programas encontrados en la máquina virtual (programa cargado, «BRAZO MIGUEL ROBADO», programas de semestres anteriores)\\
máquina virtual «WINDOWS XP» & STEP 7-Micro/WIN; carpeta compartida \codigo{\textbackslash\textbackslash vboxsvr\textbackslash MIS\_COSAS}\\
\bottomrule
\end{tabularx}
\end{center}
Los datos propios de cada PC (\codigo{usuarios.json}, \codigo{tls\_*.pem}, \codigo{registro\_gemelo.log}) están en \codigo{.gitignore}: nunca se suben.

\section{Flujo de trabajo con Git}
\begin{figure}[H]
\centering
\begin{tikzpicture}[x=1cm,y=1cm]
\draw[very thick, naranja] (0,0) -- (13,0) node[right, font=\scriptsize\ttfamily]{main};
\foreach \x in {0.5,2,6.5,8,11.5} \fill[naranja] (\x,0) circle (0.12);
\draw[thick, azulg] (2,0) .. controls (2.5,1) .. (3.2,1) -- (5.6,1) .. controls (6.2,1) .. (6.5,0);
\foreach \x in {3.2,4.4,5.6} \fill[azulg] (\x,1) circle (0.1);
\node[font=\scriptsize\ttfamily, above] at (4.4,1.1) {feature/ladder-s7200-homing};
\draw[thick, verde] (8,0) .. controls (8.5,-1) .. (9.2,-1) -- (10.6,-1) .. controls (11.2,-1) .. (11.5,0);
\foreach \x in {9.2,10.6} \fill[verde] (\x,-1) circle (0.1);
\node[font=\scriptsize\ttfamily, below] at (9.9,-1.1) {feature/db-y-supervision};
\node[etiq] at (6.5,0.45) {pull request\\revisado};
\node[etiq] at (13,-0.6) {push $\Rightarrow$ página publicada};
\end{tikzpicture}
\caption{Ramas de trabajo que se integran a \texttt{main} con revisión.}
\end{figure}

\begin{lstlisting}
git checkout -b feature/planos-v2-kaizen         # rama aislada para un componente
git add semestres/2026_II/planos/electrico/brazo_electrico_v2.dwg
git commit -m "Plano electrico v2.0: DIR+ de X e Y y garra corregidos"
git push origin feature/planos-v2-kaizen           # luego: Pull Request en GitHub
git log --oneline --graph                          # historial para el informe
\end{lstlisting}

\begin{center}\small
\begin{tabular}{@{}ll@{}}
\toprule
\textbf{Dato del historial} & \textbf{Valor (10/11/2026)}\\
\midrule
commits en \codigo{main} & 49 y los de este manual\\
integrantes con commits & 5 (Gabriel Gómez, José Lenin Mamani, Román Machicado, Eduardo Chura, \dots)\\
ramas de trabajo & \codigo{feature/ladder-s7200-homing}, \codigo{feature/db-y-supervision}, \codigo{feature/planos-v2-kaizen}, \codigo{feat/cinematica-modelo-3d}\\
primer commit & 28/09/2026\\
\bottomrule
\end{tabular}
\end{center}

\section{Convenciones}
\begin{itemize}
\item \textbf{Programas del PLC}: un archivo por versión con el número en el nombre (\codigo{brazo\_secuencias\_original\_v1.awl}); exportar siempre también el AWL para que el código se pueda leer y comparar.
\item \textbf{Mensajes de commit}: en castellano, en presente, diciendo qué cambia y por qué.
\item \textbf{Fuentes y compilados}: los documentos se guardan en LaTeX con su PDF; los archivos auxiliares de compilación están en \codigo{.gitignore}.
\item \textbf{Medidas}: cada dato lleva su origen (leída, estimada, SUPUESTO).
\end{itemize}
